\section{Cyber capability is not agent compromise}\label{sec:C28:cybercapability}
Cybersecurity capability benchmarks evaluate whether an agent can perform technical security tasks. They differ from benchmarks asking whether an attacker can redirect a benign agent. Confusing the two can produce misleading safety conclusions. A capable vulnerability analyst can be well controlled, while a weaker model can still violate confidentiality through an overprivileged connector.

CyberGym evaluates vulnerability reproduction in realistic software projects~\cite{corpus250602548}. A successful proof-of-concept reproduction demonstrates a particular capability under the benchmark's information and environment conditions. It does not establish full exploitation, remediation, or unrestricted offensive effectiveness. CyberGym-E2E broadens the evaluation toward discovery, reproduction and patch generation~\cite{corpus260604460}. The stages remain distinct because success in one does not imply success in the others.

ExploitGym studies turning known vulnerability-triggering inputs into concrete exploitation outcomes under controlled configurations~\cite{corpus260511086}. Its trusted-access and safeguard conditions matter when interpreting frontier capability. Results obtained with safeguards disabled are capability-boundary evidence, not estimates of ordinary product behavior or real-world attack prevalence. Security protections and environment variants should remain part of the reported configuration.

Patch verification also needs dual criteria: removal of the vulnerability and preservation of required functionality. A patch can eliminate the triggering behavior by disabling the program or rejecting all inputs. That is not necessarily a useful remediation. Test security and regression behavior together, and retain uncertainty where the available tests do not exhaust either property. Reference patches can supply evidence without becoming infallible specifications.

For comparative studies, match initial information, tool access, time, and verifier feedback. An agent given a vulnerability description and triggering input has a different task from one asked to discover a defect from scratch. A benchmark using real historical vulnerabilities may also have contamination concerns. Report task lineage and exposure controls rather than infer novelty from the date of the benchmark paper alone.

\subsection{Symbolic search, fuzzing, and static analysis have different oracles}\label{sec:C28:lineage}
The 2026 CyberGym and ExploitGym results belong to an older automated-security lineage, but do not supersede every mechanism in it. AEG, published at NDSS 2011, formulates control-flow exploit generation as a verification problem and uses preconditioned symbolic execution. MAYHEM, published at IEEE Security and Privacy 2012, works at the binary level and combines online and offline symbolic execution with symbolic memory-index reasoning~\cite{avgerinos2011aeg,cha2012mayhem}. These are foundational designs with historical target and mitigation assumptions. Their feasibility results should not be presented as current success rates against modern software.

The methodological contrast is important. Symbolic execution constructs path conditions and asks a solver for witnesses within a model of program semantics. Fuzzing samples or mutates inputs and retains feedback that guides further search. An agent can orchestrate either process, interpret a failure, or propose a patch; its language-model reasoning does not replace the solver's satisfiability condition or the runtime's observed failure. A benchmark must identify which stage was automated and what evidence establishes the claimed outcome. A reproducible crash, a sanitizer diagnosis, control of an execution effect, and a validated repair are different predicates.

V8's inspected JS-Fuzzer source provides a concrete continuing implementation: it consumes a corpus, generates JavaScript cases, executes shells, and supports differential testing and failure minimization~\cite{v8jsfuzzer2026}. Its current documentation also describes a local crash runner. A differential mismatch need not be a security vulnerability: nondeterminism, unsupported flags, or permitted semantic variation can produce disagreement. Deduplicate by a justified root-cause criterion rather than counting every crashing input as a distinct defect. Compare repeated campaigns with equal CPU-time envelopes and include non-discoveries; a successful anecdote gives no estimate of search efficiency.

Fault-model alignment supplies a separate axis. SbxBrk's CCS 2025 study targets V8 software-based fault isolation by mutating data as trusted code loads it from the untrusted heap~\cite{bars2025sfi}. This deliberately assumes prior control inside the sandbox; ordinary front-end input fuzzing tests a different boundary. A correct conclusion is therefore conditional isolation failure under that fault model, not a complete end-to-end compromise probability. In an agent sandbox, the analogous research question is whether trusted wrappers remain safe when an inner computation is already adversarial. Prompt obedience cannot certify that boundary.

TYPEPULSE, in USENIX Security 2025, instead analyzes Rust type conversions, trait bounds and pointer aliases, constructs a property graph, and checks patterns for misalignment, inconsistent layout and mismatched scope~\cite{chen2025typepulse}. It is a static-analysis mechanism, not an LLM agent benchmark or a guarantee for all unsafe Rust. Findings require confirmation under the analyzer's abstraction and affected program version; absence of a finding is not a completeness certificate. The cross-layer deduction is to retain distinct evidence channels: static candidate, dynamic reproduction, effect attribution, patch test, and containment result. An agent should advance a claim only as far as its strongest validated channel supports.

